\documentclass[11pt, letterpaper]{article}
\usepackage{listings}
\usepackage{xcolor}
\usepackage{array}
\usepackage{longtable}
\usepackage{parskip}
\usepackage{float}
\usepackage{graphicx}
\usepackage{hyperref}
\usepackage{xurl}
\usepackage[dvipsnames]{xcolor}
\usepackage[a4paper,top=2.54cm,bottom=2.54cm,left=3.17cm,right=3.17cm,marginparwidth=1.75cm]{geometry}

\usepackage[english]{babel}
\usepackage[utf8x]{inputenc}
\usepackage[T1]{fontenc}

\usepackage{textcomp} %plusminus in text
\usepackage[font=footnotesize]{caption} %makes captions small
\usepackage{tabularx} % extra features for tabular environment
\usepackage{multirow} 
\usepackage{amsmath}  % improve math presentation
\usepackage{siunitx}
\usepackage{hyperref}

\usepackage{titlesec} % Customising section titles
\titleformat{\section}{\large\bfseries}{\thesection.}{1em}{} % Makes section headers bold, and adds a period after number
\titleformat{\subsection}{\normalsize\bfseries}{\thesubsection.}{1em}{} % Makes section headers bold, and adds a period after number
\setlength{\parindent}{0pt}

%New colors defined below
\definecolor{codegreen}{rgb}{0,0.6,0}
\definecolor{codegray}{rgb}{0.5,0.5,0.5}
\definecolor{codepurple}{rgb}{0.58,0,0.82}
\definecolor{backcolour}{rgb}{0.95,0.95,0.92}

%Code listing style named "mystyle"
\lstdefinestyle{mystyle}{
  backgroundcolor=\color{backcolour},   commentstyle=\color{codegreen},
  keywordstyle=\color{magenta},
  numberstyle=\tiny\color{codegray},
  stringstyle=\color{codepurple},
  basicstyle=\ttfamily,
  breakatwhitespace=false,         
  breaklines=true,                 
  captionpos=b,                    
  keepspaces=true,                 
  numbers=left,                    
  numbersep=5pt,                  
  showspaces=false,                
  showstringspaces=false,
  showtabs=false,                  
  tabsize=2
}

% Windows CMD style named "cmdwinstyle"
\lstdefinestyle{cmdwinstyle}{
  backgroundcolor=\color{white},        % White background (like CMD window)
  basicstyle=\ttfamily\color{black},    % Black monospaced text
  keywordstyle=\color{blue},            % Commands in blue
  commentstyle=\color{green!60!black},  % Comments in dark green
  stringstyle=\color{red!80!black},     % Strings, if any
  numberstyle=\tiny\color{gray},        % Line numbers
  breaklines=true,
  captionpos=b,
  numbers=left,
  numbersep=5pt,
  showstringspaces=false,
  frame=single,                          % Single line frame
  rulecolor=\color{gray},                % Frame color
  tabsize=2
}

%"mystyle" code listing set
\lstset{style=mystyle}

\title{\line(1,0){250}\\A Title\\\line(1,0){250}}
\author{Kuba Kane}
\date{April 2026}


\begin{document}

\begin{titlepage}
    \begin{center}
        \vspace*{1.5cm}

        \normalsize \textbf{PHYS30762 Object Oriented Programming in C++}
        \vspace{0.5cm}
        \linebreak \Large \textbf{Analogue Circuits Report}
        \vspace{2cm}
        \linebreak \normalsize \textit{} 
        \linebreak \textit{Kuba Kane}
	\linebreak \textit{11019699}
        \vspace{12pt}
        \vspace{12pt}
        \linebreak Department of Physics and Astronomy
        \linebreak The University of Manchester
        \vspace{12pt}
        \vspace{12pt}
        \linebreak May 2026
        \vspace{\fill}
    \end{center}
    \vspace{3cm}

	\textbf{Abstract}
    
	\vspace{12pt} Abstract text goes here
    
    \vspace{\fill}
    
\end{titlepage}

\newpage
\vspace{5pt}
\tableofcontents

\newpage
\vspace{5pt}

\section{Introduction}
This program allows users to build analogue circuits using resistors, capacitors, and inductors (which are all classes derived from an abstract  base \verb|Component| class). A \verb|Circuit| class can be set to be series or parallel, and any number of components can be added to it, which is then stores as a vector of components (specifically a vector of unique pointers to components). Since the \verb|Circuit| class is inherited from the \verb|Component| class itself, it is possible to add series circuits to parallel circuits, or vice versa, and therefore construct circuits of any topology.

In the main file of this program, a \textit{library} of components is created, which stores different types of resistors, capacitors and inductors. These components are validated through the use of validation files, such that if the user tries to initialise a component with properties which do not match those specified in the validation file, an error message will be returned. The validation files specify a component name, like \textit{Wirewound} for resistor, as well as a tolerance to dissipative effects, and a maximum frequency to which they can withstand \cite{eepower_resistormaterials}. Each validation file then contains maximum and minimum possible values for the resistance, capacitance and inductance respectively. 

Due to the use of a vector of unique pointers to store the components in the \verb|Circuit| class, adding components to circuits or moving them is not trivial. The process is best-described using an analogy, considering the user of the program to be a student in an undergraduate circuits lab. In which case, the user wants to build a circuit using components from the lab's demonstrator, who has this library of components available. This library must contain components which are manufacturable, and the validation files could be seen as catalogues of different types of each component which can be manufactured.

You could think of the validation files as catalogues of components which the lab demonstrator can browse and use to choose certain materials or properties which they want to "buy" and have available in their library.


\section{Code design and implementation}
Text goes here

\section{Results}
Text goes here

\section{Discussion and conclusions}
Text goes here

\newpage
\nocite{*}
\phantomsection     
\addcontentsline{toc}{section}{References} 
\bibliographystyle{plain}  % or alpha, abbrv, unsrt, etc.
\bibliography{ref.bib} % name of your .bib file (without .bib extension)



\end{document}